% Short literature review for the MSc thesis "Where the Shortcut Sits" (October 2026).
% Uses the same bibliography as the full review (refs.bib). Build:
%   pdflatex literature_review_short && bibtex literature_review_short && pdflatex literature_review_short && pdflatex literature_review_short
\documentclass[11pt,a4paper]{article}
\usepackage[margin=2.4cm]{geometry}
\usepackage[T1]{fontenc}
\usepackage{lmodern}
\usepackage{booktabs,array}
\usepackage[round,authoryear]{natbib}
\usepackage[hidelinks]{hyperref}
\setlength{\parskip}{4pt}
\setlength{\emergencystretch}{2em}
\newcolumntype{P}[1]{>{\raggedright\arraybackslash}p{#1}}

\title{Image artifacts, region-of-interest masking and shortcut learning in medical image classification: a review
of work published from 2021 to 2026}
\author{MSc thesis \emph{Where the Shortcut Sits} --- literature review}
\date{October 2026}

\begin{document}
\maketitle

\section{Introduction}
A classifier takes a shortcut when it relies on something that predicts the label in the training data without being
evidence of the disease \citep{geirhos2020shortcut}. In medical imaging such signals are rarely random. Clinicians
measure suspicious nodules, mark lesions they intend to excise and drain collapsed lungs, so calipers, skin markings
and chest tubes appear more often in the images of sick patients. A common response is to segment the lesion or organ
and let the model see only that region of interest (ROI). The thesis asks when this works, and the answer it gives is
that it depends on where the artifact lies: masking removes an artifact outside the ROI but keeps one inside it, and the
model then relies on it more.

This review covers peer-reviewed work from 2021 to October 2026 that bears on that question. We kept to journals and
main-track conferences with an established reputation in medical imaging and machine learning, and left out preprints
and workshop papers; a few older papers are cited where later work builds on them. The review is organised in seven
parts: the evidence that medical classifiers take shortcuts, the artifacts documented in each modality, what is known
about restricting a model to the ROI, how shortcuts are detected, how they are removed, why models prefer them, and how
claims of this kind should be evaluated. A final section states where the thesis stands relative to this work.

\section{Evidence that medical classifiers take shortcuts}
The COVID-19 literature made the problem hard to ignore. \citet{degrave2021ai} showed that radiographic COVID-19
detectors relied on laterality markers, patient positioning and other differences between the source datasets, and
that their accuracy fell on data from other hospitals; they used saliency maps and generative counterfactuals to show
what the models attended to. A systematic review of 62 published COVID-19 imaging models found none fit for clinical
use, mainly because of bias in small datasets and the variability of large ones assembled from many sources
\citep{roberts2021pitfalls}. The most direct demonstration came from \citet{maguolo2021critic}, who still classified
COVID-19 radiographs after blacking out the lungs. If the label can be predicted with the anatomy removed, the
information must be somewhere else in the image. This is exactly the reasoning behind ROI masking, and it is worth
noting that it says nothing about shortcuts that sit \emph{inside} the anatomy.

A second line of work concerns patient attributes. Deep models predict self-reported race from medical images, even
when the images are heavily degraded \citep{gichoya2022race}, and chest-radiograph classifiers underdiagnose
under-served groups, most strongly at intersections such as Hispanic female patients
\citep{seyyedkalantari2021underdiagnosis}. \citet{glocker2023encoding} found that disease models encode protected
characteristics in their features, and \citet{yang2024limits} showed that the models encoding demographics most
strongly have the largest fairness gaps, and that a model made fair in one hospital is not necessarily fair in another.
Encoding an attribute is not the same as using it, and \citet{brown2023shortcut} proposed shortcut testing to
distinguish the two by varying how strongly the attribute is encoded. \citet{jones2024causal} argue that dataset biases
which look alike in the data can arise from different causal mechanisms (prevalence, presentation and annotation) and
need different remedies. We find this framing useful: the artifacts in the thesis are produced by clinical actions
taken on suspicion, so they are part of how the data are generated and cannot be cleaned away as noise.

Several groups study shortcuts under controlled conditions. \citet{stanley2024simba} built synthetic brain images with
known disease and bias effects, compared three mitigation strategies and found reweighing to work best;
\citet{sun2023right} added a tag, hyperintensities or an obstruction to chest radiographs and found that explanation
methods differ widely in whether they reveal the injected signal. These designs fix where the bias is placed. None of
them varies its position relative to the region that matters, which is the variable the thesis manipulates.

\section{Artifacts documented in practice}
Dermoscopy has the richest record. Ruler scale bars changed the melanoma predictions of a market-approved classifier
\citep{winkler2021scalebars}, as surgical skin markings had done in an earlier study \citep{winkler2019markings}.
Explanation-based audits of ISIC models uncovered colour calibration patches, rulers and other artifacts
\citep{anders2022clever,pahde2023r2r}. Hair that covers the lesion degrades classifiers \citep{wang2024hair}, and for
dark corners, adding synthetic ones during training helped more than removing the real ones \citep{pewton2024dca}. The
ISIC collections also contain duplicate images across releases, which can leak between training and test sets
\citep{cassidy2022isic}. Read together, these studies show that the artifacts sit in different places: patches and dark
corners lie mostly outside the lesion, rulers near the image edge, markings around the lesion, and hair crosses it. The
studies report the effect of an artifact being present; none asks whether its position matters.

In ultrasound the evidence is thinner. \citet{lin2024segshortcut} showed that calipers and on-screen text act as
shortcuts in fetal ultrasound, even for segmentation. We found no peer-reviewed study that treats calipers as a
classification shortcut in thyroid or ovarian ultrasound, although calipers are placed on the nodule border by design.
In chest radiography, pneumothorax classifiers respond to the chest tube that treats the condition.
\citet{rueckel2021pneumothorax} reduced this confounding by training with in-image annotations, and
\citet{jimenez2023shortcuts} labelled drains on a public dataset to expose the shortcut. A drain projects over the lung
field, so a lung mask cannot remove it. For capsule endoscopy we found no study of bubbles or debris as a shortcut, even
though debris and lesions often occupy the same pixels.

\section{Does restricting the model to the ROI help?}
In natural images, classifiers can recognise objects from the background alone, and more accurate models tend to rely
on it less \citep{xiao2021noise}. The same holds for medical images. In a careful masking study on chest radiographs
and fundus photographs, models trained on full images performed better on images with the ROI removed than on images
containing only the ROI \citep{sourget2025mask}. Training on segmented dermoscopy images improved balanced accuracy for
classifiers trained on HAM10000 but not for those trained on ISIC \citep{maron2021confounding}, a mixed result that is
easy to overlook. ISNet takes a softer route and penalises relevance outside the segmentation mask during training,
which made chest-radiograph classifiers more robust to background bias \citep{bassi2024isnet}.

All of these studies assume that the shortcut lies in the background. The one result that points the other way comes
from natural images: when several shortcuts are present, removing one can make the model rely more on the others
\citep{li2023whacamole}. ROI masking is a mitigation of exactly this kind. If an artifact remains inside the ROI, the
same reasoning predicts that masking will increase reliance on it. To our knowledge this has not been tested in medical
images, and it is the starting point of the thesis.

\section{Detecting shortcuts}
Explanation methods are the most common tool. Spectral relevance analysis clusters heatmaps to reveal Clever-Hans
strategies \citep{anders2022clever}, and Reveal-to-Revise turns detection, correction and re-evaluation into a cycle
\citep{pahde2023r2r}. The limits are well documented: saliency maps localise chest pathology worse than radiologists
\citep{saporta2022saliency}, and their ability to reveal a known confounder varies from method to method
\citep{sun2023right}. In our view there is a further limit for in-ROI artifacts: when calipers sit on a nodule, a heatmap
that highlights the nodule cannot tell whether the model looks at the tissue or at the calipers.

Counterfactual methods avoid this by changing the image. Diffusion-based counterfactuals add or remove a suspected
shortcut and measure the change in prediction \citep{weng2024fastdime}, and RadEdit edits chest radiographs under masks
to stress-test classifiers \citep{perezgarcia2024radedit}. The thesis uses the same idea in its simplest form, by
pasting one real artifact inside or outside the ROI of the same images.

\section{Removing shortcuts}
With group labels, retraining only the last layer on balanced data recovers much of the robustness lost to spurious
correlations \citep{kirichenko2023dfr}, because standard training already learns the core features even when the
classifier does not use them \citep{izmailov2022feature}. Without group labels, the errors of a first model can be
up-weighted \citep{liu2021jtt}, or the last layer reweighted automatically \citep{qiu2023afr}; variance-of-risk
penalties are another option \citep{krueger2021rex}. The benchmarks are sobering. Across 20 algorithms and 12 datasets,
methods improved subgroup robustness only for some kinds of shift \citep{yang2023change}, and in medical imaging,
bias-mitigation methods did not improve fairness significantly over standard training, while the choice of model
selection mattered more \citep{zong2023medfair}.

A more principled line starts from causal reasoning. Counterfactual invariance requires that the prediction does not
change when only the nuisance changes; when the label causes the image, as in diagnosis, this implies independence of
the prediction and the nuisance given the label \citep{veitch2021counterfactual}. \citet{makar2022causally} enforce such
conditional independences with auxiliary labels available at training time, and NuRD trains under a distribution in
which the nuisance--label relation is broken \citep{puli2022nurd}. \citet{kaur2023modeling} show that the right
constraint depends on how the data were generated, so no single constraint suits every shift. Linear concept erasure
removes a concept from the features in closed form \citep{belrose2023leace}, at the cost of also removing label signal
when the two are correlated; SPLINCE was designed to keep that signal \citep{holstege2025splince}. Generative data are
a different route: diffusion-generated images improved fairness out of distribution in histopathology, chest
radiography and dermatology \citep{ktena2024generative}.

The remedy the thesis recommends, a class-conditional mean constraint on the masked features, is a first-moment, linear
case of the conditional independence of \citet{veitch2021counterfactual} and \citet{makar2022causally}. What the thesis
adds is the reason for applying it after masking, and a registered test of whether any remedy beats masking across all
cohorts. Given the benchmark results above, it is not surprising that none does.

\section{Why models prefer shortcuts}
Theory explains why a simple feature wins even when a reliable one is available. \citet{nagarajan2021understanding}
show that geometric and statistical skews push standard training toward spurious features even when the invariant
feature predicts the label perfectly. Gradient descent with cross-entropy can starve less salient features
\citep{pezeshki2021gradient}, and its implicit preference for large margins produces shortcut reliance even when the
shortcut adds no information \citep{puli2023dontblame}. \citet{hermann2024foundations} separate a feature's
\emph{predictivity} from its \emph{availability}, that is, how easily it can be extracted, and measure shortcut bias as
reliance on the more available but less predictive feature. For frozen pretrained features, as in most of the thesis,
linear probing can generalise better out of distribution than full fine-tuning \citep{kumar2022finetuning}. The
thesis's linear-Gaussian account fits this picture: masking removes competing evidence outside the ROI and so changes
the relative predictivity of what is left. The difference is that the account was used to make quantitative predictions
that were registered before the results existed.

\section{How such claims should be evaluated}
\citet{varoquaux2022failures} review methodological failures in medical-imaging machine learning, from biased data to
publication incentives, and \citet{kapoor2023leakage} found data leakage, including non-independent splits, in at least
294 papers across 17 fields. Results also vary with data sampling, initialisation and hyperparameters, and accounting
for several of these sources at once gives better estimates of that variation \citep{bouthillier2021variance}. In MICCAI
papers, confidence intervals are rarely reported, and their typical width is larger than the gap between the first- and
second-ranked methods \citep{christodoulou2024ci}. Overall metrics also hide failures in clinically important subgroups
\citep{yang2024limits}, a problem first described as hidden stratification \citep{oakden2020hidden}. These papers shaped
the evaluation in the thesis: splits grouped by patient or lesion, a bootstrap that resamples training seeds and test
images together, registration before results, and evaluation on the cases that contradict the shortcut as well as on
the whole test set.

\section{Where the thesis stands}
The literature establishes four things firmly. Medical classifiers take shortcuts, including ones produced by clinical
practice; the artifacts involved are well documented in dermoscopy and chest radiography; image backgrounds carry
signal, and restricting a model to the ROI sometimes helps; and no mitigation method is reliably better than standard
training across settings. Table~\ref{tab:gaps} lists what remained open and how the thesis addresses it.

\begin{table}[ht]\centering\small
\caption{Open questions in the reviewed work and the corresponding contribution of the thesis.}\label{tab:gaps}
\begin{tabular}{P{4.6cm}P{4.6cm}P{5.6cm}}\toprule
Open question & Closest prior work & Thesis \\\midrule
Does the artifact's location relative to the ROI decide whether masking works? & \citet{sourget2025mask,maron2021confounding,bassi2024isnet}: masking studies that assume a background shortcut & Location varied directly: controlled sweeps, traps built from real artifacts, and a transplant of the same artifact across the ROI boundary \\
Can masking increase reliance on an artifact? & \citet{li2023whacamole}: removing one shortcut amplifies others (natural images) & Reliance on an in-ROI artifact grows after masking, in four cohorts and three modalities \\
Can the harm be predicted? & \citet{hermann2024foundations,puli2023dontblame}: explanatory theories & A two-parameter account with predictions registered before the results \\
What should be used instead of masking? & \citet{veitch2021counterfactual,makar2022causally,zong2023medfair} & A registered test against masking; a class-conditional constraint comes closest, none dominates \\
Calipers and debris as classification shortcuts & \citet{lin2024segshortcut} (segmentation only) & Thyroid and ovarian ultrasound and capsule endoscopy alongside dermoscopy \\\bottomrule
\end{tabular}\end{table}

Some of this is not new, and the thesis should say so. That a mask cannot remove what it keeps is self-evident; what was
missing is evidence that masking makes reliance on the kept artifact grow, how large the effect is on real artifacts,
and how to predict it. The recommended constraint is a special case of known conditional-invariance methods, and the
group-robustness methods are used as baselines rather than proposed. The closest conceptual precedent is the
whac-a-mole result of \citet{li2023whacamole}; the thesis supplies the location-specific mechanism and the medical
evidence.

\section*{Use of AI tools}
This review was prepared with the assistance of an AI system (Claude, Anthropic), which was used to search the
literature, check bibliographic details against publisher and proceedings pages, and draft the text. The choice of
scope, the interpretation of the studies and any remaining errors are the author's responsibility.

\bibliographystyle{plainnat}
\bibliography{refs}
\end{document}
